% !TeX encoding = UTF-8
% !TEX program = xelatex
\documentclass[12pt,a4paper]{ctexart}

% === Packages ===
\usepackage{amsmath,amssymb,amsthm}
\usepackage{mathtools}
\usepackage{geometry}
\usepackage{tikz}
\usetikzlibrary{arrows.meta, positioning, shapes.geometric, shadows, fit, backgrounds}
\usepackage{hyperref}
\usepackage{booktabs}
\usepackage{tabularx}
\usepackage{enumitem}
\usepackage{bm}
\usepackage{tcolorbox}
\usepackage{cleveref}
\usepackage{fancyhdr}
\usepackage{float}

\geometry{a4paper, margin=2.5cm, headheight=15pt}
\hypersetup{hidelinks}

% === TOC: no indent for any level ===
\makeatletter
\renewcommand*\l@section{\@dottedtocline{1}{0em}{1.8em}}
\renewcommand*\l@subsection{\@dottedtocline{2}{0em}{2.4em}}
\renewcommand*\l@subsubsection{\@dottedtocline{3}{0em}{3.0em}}
\makeatother

% === Theorem environments ===
\newtheorem{theorem}{定理}[section]
\newtheorem{lemma}[theorem]{引理}
\newtheorem{corollary}[theorem]{推论}
\newtheorem{proposition}[theorem]{命题}
\theoremstyle{definition}
\newtheorem{definition}[theorem]{定义}
\newtheorem{remark}[theorem]{注记}
\theoremstyle{definition}
\newtheorem{example}[theorem]{例}

% === Custom commands ===
\newcommand{\vect}[1]{\bm{#1}}
\newcommand{\mat}[1]{\bm{#1}}
\newcommand{\trans}{^{\mathsf{T}}}
\newcommand{\inv}{^{-1}}
\newcommand{\norm}[1]{\lVert#1\rVert}
\newcommand{\abs}[1]{\lvert#1\rvert}
\newcommand{\inner}[2]{\langle#1,#2\rangle}
\newcommand{\simto}{\sim}
\newcommand{\real}{\mathbb{R}}
\newcommand{\proj}{\mathbb{P}}
\DeclareMathOperator*{\argmin}{argmin}
\DeclareMathOperator*{\argmax}{argmax}

% === Result box ===
\newtcolorbox{resultbox}[1][]{
  colback=blue!5!white,
  colframe=blue!75!black,
  fonttitle=\bfseries,
  title={#1},
  arc=4pt,
  before skip=12pt,
  after skip=12pt
}

\newtcolorbox{finalbox}{
  colback=red!3!white,
  colframe=red!70!black,
  fonttitle=\bfseries,
  title={本章最终LUT公式},
  arc=4pt,
  before skip=12pt,
  after skip=12pt
}

% === Page style ===
\pagestyle{fancy}
\fancyhf{}
\fancyhead[L]{\small 基于预标定与查表的全景拼接——数学推理}
\fancyhead[R]{\small\thepage}
\renewcommand{\headrulewidth}{0.4pt}

% ============================================
\title{\textbf{基于预标定与查表的图像全景拼接算法\\[6pt] 数学推理与LUT变换公式}}
\author{K. Lee}
\date{\today}

\begin{document}

\maketitle
\thispagestyle{empty}

\begin{abstract}
\noindent 本文系统推导了基于相机预标定参数与离线查表（Look-Up Table, LUT）进行图像全景拼接的完整数学框架。
涵盖\emph{相机标定方法}（张氏标定、立体标定、多相机阵列标定、误差传播分析），
按\emph{变换方式}（纯旋转、纯平移、旋转+平移）、\emph{相机数量}（双相机、多相机阵列）、\emph{光心是否重合}、\emph{有无畸变}四个维度，
分章节给出每种情形下的精确或近似LUT变换矩阵/公式。
每章末尾以加框公式给出该情形对应的最终查表映射，可直接用于工程实现。
\end{abstract}

\newpage
\pagestyle{empty}
\tableofcontents

\newpage
\pagenumbering{arabic}
\setcounter{page}{1}
\pagestyle{fancy}

% ============================================================
\include{chapters/fundamentals}
% ============================================================
% Chapter: 相机标定方法
% ============================================================
\section{相机标定方法}
\label{sec:calibration}

前文已建立LUT拼接的完整数学框架，其全部依赖已知的标定参数——内参矩阵$\mat{K}_i$、畸变系数$(k_1,k_2,k_3,p_1,p_2)$、相机间外参$(\mat{R}_{ij},\vect{T}_{ij})$以及世界坐标系到各相机的旋转$\mat{R}_{iw}$。本章系统推导这些参数的标定方法及标定误差对LUT的影响。

\subsection{单相机内参标定——张正友平面标定法}
\label{sec:zhang}

张氏标定法（Zhang, 2000）是目前最广泛使用的相机标定方法，仅需拍摄若干不同姿态下的平面棋盘格图案。

\subsubsection{基本模型}

设平面标定板位于世界坐标系$Z_w=0$平面，其上的三维点$\vect{X}_w = [X_w, Y_w, 0]\trans$。相机投影方程为：
\begin{equation}
\vect{p} \simto \mat{K}\bigl[\mat{R} \mid \vect{t}\bigr]
\begin{bmatrix}
X_w \\ Y_w \\ 0 \\ 1
\end{bmatrix}
= \mat{K}\bigl[\vect{r}_1 \;\; \vect{r}_2 \;\; \vect{t}\bigr]
\begin{bmatrix}
X_w \\ Y_w \\ 1
\end{bmatrix}
\label{eq:zhang-proj}
\end{equation}
其中$\mat{R} = [\vect{r}_1\; \vect{r}_2\; \vect{r}_3]$，$\vect{t}$为世界坐标系到相机坐标系的变换。

由于$Z_w=0$，三维到二维的映射退化为平面单应：
\begin{equation}
\vect{p} \simto \mat{H}
\begin{bmatrix}
X_w \\ Y_w \\ 1
\end{bmatrix},
\qquad
\mat{H} = \mat{K}\bigl[\vect{r}_1 \;\; \vect{r}_2 \;\; \vect{t}\bigr]
\label{eq:zhang-H}
\end{equation}

\subsubsection{单应矩阵估计}

对每幅标定图像$j$（$j = 1, \dots, M$），从$N$个角点（$N \ge 4$）的像素坐标$(u_{j,n}, v_{j,n})$与世界坐标$(X_{w,n}, Y_{w,n})$的对应关系，通过直接线性变换（DLT）估计$3\times 3$单应矩阵$\mat{H}_j$。

对每个角点$n$，齐次映射$\vect{p}_n \simto \mat{H}\,\vect{X}_{w,n}$写出两个线性方程（消去齐次因子）：
\begin{align}
\vect{h}_1\trans\vect{X}_{w,n} - u_n\,\vect{h}_3\trans\vect{X}_{w,n} &= 0 \label{eq:dlt-1} \\
\vect{h}_2\trans\vect{X}_{w,n} - v_n\,\vect{h}_3\trans\vect{X}_{w,n} &= 0 \label{eq:dlt-2}
\end{align}
其中$\vect{h}_k\trans$为$\mat{H}$的第$k$行。$N$个点提供$2N$个方程，可用SVD求解（$\mat{H}$有8个自由度，需$N \ge 4$）。

\subsubsection{解析求解内参——旋转矩阵的正交约束}

单应矩阵$\mat{H} = [\vect{h}_1\; \vect{h}_2\; \vect{h}_3]$已估计。由\eqref{eq:zhang-H}：
\begin{equation}
[\vect{h}_1 \;\; \vect{h}_2 \;\; \vect{h}_3] = \lambda\,\mat{K}\,[\vect{r}_1 \;\; \vect{r}_2 \;\; \vect{t}]
\label{eq:H-decomp}
\end{equation}
其中$\lambda$为任意尺度因子（$\mat{H}$只能确定到尺度）。

利用$\vect{r}_1,\vect{r}_2$的正交性（$\mat{R}$为旋转矩阵）：
\begin{align}
\vect{r}_1\trans\vect{r}_2 &= 0 \label{eq:orth1} \\
\vect{r}_1\trans\vect{r}_1 = \vect{r}_2\trans\vect{r}_2 &= 1 \label{eq:orth2}
\end{align}

由\eqref{eq:H-decomp}：$\vect{r}_1 = \lambda\inv\mat{K}\inv\vect{h}_1,\; \vect{r}_2 = \lambda\inv\mat{K}\inv\vect{h}_2$。代入\eqref{eq:orth1}--\eqref{eq:orth2}得到两个对内参$\mat{K}$的约束：
\begin{align}
\vect{h}_1\trans\mat{K}^{-\mathsf{T}}\mat{K}\inv\vect{h}_2 &= 0 \label{eq:con1} \\
\vect{h}_1\trans\mat{K}^{-\mathsf{T}}\mat{K}\inv\vect{h}_1 &=
\vect{h}_2\trans\mat{K}^{-\mathsf{T}}\mat{K}\inv\vect{h}_2 \label{eq:con2}
\end{align}

\begin{definition}[绝对二次曲线的像（Image of Absolute Conic, IAC）]
定义对称矩阵$\mat{B} = \mat{K}^{-\mathsf{T}}\mat{K}\inv$。展开后$\mat{B}$由6个独立参数表示：
\begin{equation}
\mat{B} =
\begin{bmatrix}
B_{11} & B_{12} & B_{13} \\
B_{12} & B_{22} & B_{23} \\
B_{13} & B_{23} & B_{33}
\end{bmatrix}
\label{eq:B-matrix}
\end{equation}
\end{definition}

每个单应矩阵提供两个约束方程\eqref{eq:con1}--\eqref{eq:con2}，可写为关于$\vect{b} = [B_{11}, B_{12}, B_{22}, B_{13}, B_{23}, B_{33}]\trans$的线性方程：
\begin{equation}
\begin{bmatrix}
\vect{v}_{12}\trans \\ (\vect{v}_{11} - \vect{v}_{22})\trans
\end{bmatrix}
\vect{b} = \vect{0}
\label{eq:v-matrix}
\end{equation}
其中$\vect{v}_{pq}$为由$\mat{H}$的列构造的$6\times 1$向量。$M \ge 3$幅不同姿态的图像可提供$2M \ge 6$个方程，通过SVD求解$\vect{b}$（进而得到$\mat{B}$）。

\subsubsection{从\texorpdfstring{$\mat{B}$}{B}解析提取内参}

由$\mat{B} = \mat{K}^{-\mathsf{T}}\mat{K}\inv$，可解析反解出各内参：
\begin{align}
c_y &= \frac{B_{12}B_{13} - B_{11}B_{23}}{B_{11}B_{22} - B_{12}^2} \label{eq:extract-cy} \\[4pt]
\lambda_0 &= B_{33} - \frac{B_{13}^2 + c_y(B_{12}B_{13} - B_{11}B_{23})}{B_{11}} \label{eq:extract-lambda} \\[4pt]
f_x &= \sqrt{\frac{\lambda_0}{B_{11}}} \label{eq:extract-fx} \\[4pt]
f_y &= \sqrt{\frac{\lambda_0 B_{11}}{B_{11}B_{22} - B_{12}^2}} \label{eq:extract-fy} \\[4pt]
s   &= -\frac{B_{12} f_x^2 f_y}{\lambda_0} \qquad (\text{skew，通常}\, \approx 0) \label{eq:extract-s} \\[4pt]
c_x &= \frac{s c_y}{f_y} - \frac{B_{13} f_x^2}{\lambda_0} \label{eq:extract-cx}
\end{align}
当假设像素无倾斜（$s=0$，现代CMOS传感器普遍满足）时，上述公式简化为：
\begin{align}
f_x &= \sqrt{\frac{\lambda_0}{B_{11}}}, \qquad
f_y = \sqrt{\frac{\lambda_0 B_{11}}{B_{11}B_{22} - B_{12}^2}} \label{eq:simple-f} \\
c_x &= -\frac{B_{13}f_x^2}{\lambda_0}, \qquad
c_y = \frac{B_{12}B_{13} - B_{11}B_{23}}{B_{11}B_{22} - B_{12}^2} \label{eq:simple-c}
\end{align}
\subsubsection{求解外参}
对内参$\mat{K}$已知的每幅图像，外参可从单应矩阵恢复（$\lambda = 1/\|\mat{K}\inv\vect{h}_1\|$）：
\begin{align}
\vect{r}_1 &= \lambda \mat{K}\inv \vect{h}_1 \label{eq:ext-r1} \\
\vect{r}_2 &= \lambda \mat{K}\inv \vect{h}_2 \label{eq:ext-r2} \\
\vect{r}_3 &= \vect{r}_1 \times \vect{r}_2 \label{eq:ext-r3} \\
\vect{t}   &= \lambda \mat{K}\inv \vect{h}_3 \label{eq:ext-t}
\end{align}
由于噪声，$[\vect{r}_1\; \vect{r}_2\; \vect{r}_3]$未必严格正交，需通过SVD投影到最近$SO(3)$矩阵：$\mat{R} = \mat{U}\,\text{diag}(1,1,\det(\mat{U}\mat{V}\trans))\,\mat{V}\trans$。
\subsubsection{非线性优化——Levenberg--Marquardt}
解析解作为初值，通过最小化重投影误差对全部参数进行非线性优化：
\begin{equation}
\min_{\mat{K},\{\mat{R}_j,\vect{t}_j\},\{k_\ell,p_m\}}
\sum_{j=1}^{M}\sum_{n=1}^{N}
\Bigl\|\vect{p}_{j,n}^{\text{(obs)}} - \hat{\vect{p}}_{j,n}\bigl(\mat{K}, \mat{R}_j, \vect{t}_j, k_1,k_2,k_3,p_1,p_2, \vect{X}_{w,n}\bigr)\Bigr\|^2
\label{eq:lm-opt}
\end{equation}
其中$\hat{\vect{p}}_{j,n}$为投影函数（含畸变模型\eqref{eq:undistort-x}--\eqref{eq:undistort-y}）。此即\textbf{张氏标定法的完整优化目标}。
\begin{resultbox}[单相机标定所需的输入与输出]
\textbf{输入：}$M$幅（建议$10\sim20$幅）不同姿态下的棋盘格图像，棋盘格角点数$N \times N$，方格物理尺寸。
\textbf{输出：}
\[
\mat{K} = \begin{bmatrix} f_x & s & c_x \\ 0 & f_y & c_y \\ 0 & 0 & 1 \end{bmatrix},\quad
\text{distCoeffs} = [k_1, k_2, p_1, p_2, k_3]\trans
\]
\end{resultbox}

% --------------------------------------------------
\subsection{双相机外参标定——立体标定}
\label{sec:stereo-calib}

对于双相机系统，需标定相机间的相对位姿$(\mat{R}_{21}, \vect{T}_{21})$。

\subsubsection{方法一：同时立体标定}

两个相机同时拍摄同一标定板的多组姿态（两相机均能看到棋盘格）。在每个姿态$j$下，利用\cref{sec:zhang}的方法分别求解各相机的外参$(\mat{R}_{1,j}, \vect{t}_{1,j})$和$(\mat{R}_{2,j}, \vect{t}_{2,j})$（均相对于世界坐标系，即标定板）。

相机2到相机1的变换为：
\begin{equation}
\mat{R}_{21} = \frac{1}{M}\sum_{j=1}^{M} \mat{R}_{1,j}\,\mat{R}_{2,j}\trans
\quad(\text{平均后投影到}SO(3))
\label{eq:R21-stereo}
\end{equation}
\begin{equation}
\vect{T}_{21} = \frac{1}{M}\sum_{j=1}^{M} \bigl(\vect{t}_{1,j} - \mat{R}_{1,j}\,\mat{R}_{2,j}\trans\,\vect{t}_{2,j}\bigr)
\label{eq:T21-stereo}
\end{equation}

推导依据：$\vect{X}_1 = \mat{R}_{1,j}\vect{X}_w + \vect{t}_{1,j}$，$\vect{X}_2 = \mat{R}_{2,j}\vect{X}_w + \vect{t}_{2,j}$。消去$\vect{X}_w$：
\[
\vect{X}_1 = \mat{R}_{1,j}\mat{R}_{2,j}\trans(\vect{X}_2 - \vect{t}_{2,j}) + \vect{t}_{1,j}
          = \underbrace{\mat{R}_{1,j}\mat{R}_{2,j}\trans}_{\mat{R}_{21}}\vect{X}_2 + \underbrace{(\vect{t}_{1,j} - \mat{R}_{1,j}\mat{R}_{2,j}\trans\vect{t}_{2,j})}_{\vect{T}_{21}}
\]

\subsubsection{方法二：已知各自内参后的本质矩阵分解}

若两相机内参已单独标定，可通过匹配特征点计算本质矩阵$\mat{E} = \mat{K}_2\trans\,\mat{F}\,\mat{K}_1$（$\mat{F}$为基础矩阵），然后对$\mat{E}$进行SVD分解恢复$(\mat{R}_{21}, \vect{T}_{21})$（有四组可能解，通过chirality check选取正确解）。此方法可在线进行，不依赖标定板同时可见。

\subsubsection{联合优化——立体Bundle Adjustment}

将两相机的内外参与畸变联合进行LM优化：
\begin{equation}
\min \sum_{j=1}^{M}\sum_{i=1}^{2}\sum_{n=1}^{N}
\bigl\|\vect{p}_{i,j,n}^{\text{(obs)}} - \hat{\vect{p}}_{i,j,n}(\mat{K}_i, \mat{R}_{21}, \vect{T}_{21}, \mat{R}_{1,j}, \vect{t}_{1,j}, \text{dist}_i)\bigr\|^2
\label{eq:stereo-ba}
\end{equation}

\begin{finalbox}
双相机外参标定结果：
\[
\mat{R}_{21} \in SO(3),\qquad \vect{T}_{21} \in \real^3
\]
此为LUT单应矩阵$\mat{H}_{21} = \mat{K}_2\mat{R}_{21}\mat{K}_1\inv$（\cref{eq:H21}）或平面诱导单应$\mat{H}_{21}^{\Pi}$（\cref{eq:H-plane-general}）的核心输入。
\end{finalbox}

% --------------------------------------------------
\subsection{多相机阵列标定}
\label{sec:multi-calib}

$N$个相机构成刚性阵列时，需确定每个相机在世界（全景）坐标系下的位姿$(\mat{R}_{iw}, \vect{T}_{i0})$。

\subsubsection{方法一：链式标定}

若相邻相机之间有足够重叠视场，可通过逐对标定构建位姿图（Pose Graph）：
\begin{enumerate}[leftmargin=*]
  \item 对所有有重叠的相机对$(i,j)$，进行双目标定得到$(\mat{R}_{ij}, \vect{T}_{ij})$；
  \item 选定参考相机0（$\mat{R}_{00}=\mat{I}, \vect{T}_{00}=\vect{0}$），通过位姿链传播：
  \[
  \mat{R}_{i0} = \mat{R}_{i,i-1}\,\mat{R}_{i-1,i-2}\,\cdots\,\mat{R}_{10},
  \qquad
  \vect{T}_{i0} = \vect{T}_{i,i-1} + \mat{R}_{i,i-1}\vect{T}_{i-1,i-2} + \cdots
  \]
  \item 通过位姿图优化（Pose Graph Optimization）消除累积误差。
\end{enumerate}

\subsubsection{方法二：全局Bundle Adjustment}

在所有相机同时观测到标定板的姿态下（或使用多个标定板覆盖各相机视场），进行全局联合优化：
\begin{equation}
\min \sum_{i=1}^{N}\sum_{j=1}^{M}\sum_{n=1}^{N_{\text{pts}}}
\bigl\|\vect{p}_{i,j,n}^{\text{(obs)}} - \hat{\vect{p}}_{i,j,n}(\mat{K}_i, \mat{R}_{i0}, \vect{T}_{i0}, \mat{R}_{w,j}, \vect{t}_{w,j}, \text{dist}_i)\bigr\|^2
\label{eq:global-ba}
\end{equation}

\subsubsection{无重叠视场的阵列标定}

当相机间无重叠视场时（如环视车载相机），可使用：
\begin{enumerate}[leftmargin=*]
  \item \textbf{大型标定场}：在标定车间地面绘制全局标定图案，各相机从各自角度观测不同区域，通过已知的全局图案坐标关联各相机；
  \item \textbf{手眼标定（Hand-Eye Calibration）}：利用相机阵列的刚性约束$\mat{R}_{ij} = \mat{R}_i\mat{R}_j\trans$，通过运动恢复结构（SfM）求解各相机位姿；
  \item \textbf{镜面/折射辅助}：使用平面镜使无重叠的相机观测到同一标定板。
\end{enumerate}

\begin{finalbox}
多相机阵列标定结果（用于LUT全景拼接）：
\[
\mat{K}_i,\quad \mat{R}_{iw} = \mat{R}_{i0}\inv,\quad \text{distCoeffs}_i \qquad (i = 0,1,\dots,N-1)
\]
其中$\mat{R}_{iw}$直接用于球面全景LUT公式$\vect{p}_i \simto \mat{K}_i\mat{R}_{iw}\vect{d}_w$（\cref{eq:multi-rot-main}），$\vect{T}_{i0}$用于平面诱导单应或含深度映射。
\end{finalbox}

% --------------------------------------------------
\subsection{标定误差对LUT精度的影响分析}
\label{sec:calib-error-propagation}

标定参数的不确定性会传播到LUT坐标。本节给出各参数误差的一阶灵敏度分析。

\subsubsection{内参误差的像素映射偏差}

设内参误差$\Delta\mat{K}_i = \mat{K}_i^{\text{(true)}} - \mat{K}_i^{\text{(est)}}$。对于任意全景方向$\vect{d}_w$，纯旋转LUT的像素偏差为：
\begin{equation}
\Delta\vect{p}_i \approx \frac{\partial \vect{p}_i}{\partial \mat{K}_i}\,\Delta\mat{K}_i
=
\frac{\partial}{\partial\mat{K}_i}\Bigl(\mat{K}_i\mat{R}_{iw}\vect{d}_w\Bigr)\,\Delta\mat{K}_i
\label{eq:k-error}
\end{equation}

展开（令$\vect{d}_i = \mat{R}_{iw}\vect{d}_w = [x,y,z]\trans$）：
\begin{align}
\Delta u_i &\approx \frac{x}{z}\,\Delta f_x + \Delta c_x
            - \frac{f_x x}{z^2}\,\Delta z \quad(\Delta z\approx 0\text{ 对旋转}) \label{eq:du-dk} \\
\Delta v_i &\approx \frac{y}{z}\,\Delta f_y + \Delta c_y \label{eq:dv-dk}
\end{align}

\begin{proposition}[内参误差界]
对于视场角FOV内的全景像素，主点误差$\Delta c_x, \Delta c_y$造成固定的像素偏移，焦距相对误差$\Delta f_x/f_x$造成的偏移与方向角$\arctan(x/z)$成正比，在图像边缘（大角度）最显著：
\[
|\Delta u_i|_{\max} \approx |\Delta c_x| + \tan(\text{FOV}/2)\cdot\frac{|\Delta f_x|}{f_x}\cdot f_x
\]
\end{proposition}

\subsubsection{旋转标定误差的像素映射偏差}

设旋转误差$\Delta\mat{R}$（微小旋转，可用旋转向量$\Delta\vect{\omega}$表示，$\|\Delta\vect{\omega}\| \ll 1$）：
\begin{equation}
\vect{p}_i^{\text{(est)}} \simto \mat{K}_i\,\exp([\Delta\vect{\omega}]_\times)\,\mat{R}_{iw}^{\text{(true)}}\,\vect{d}_w
\label{eq:r-error}
\end{equation}

一阶近似下（令$\mat{R}_{iw}^{\text{(true)}}\vect{d}_w = [x,y,z]\trans$）：
\begin{align}
\Delta u_i &\approx \frac{f_x}{z^2}\Bigl(
              -\Delta\omega_z\,x\,y + \Delta\omega_y\,(x^2+z^2) - \Delta\omega_x\,y\,z
              \Bigr) \label{eq:du-dr} \\
\Delta v_i &\approx \frac{f_y}{z^2}\Bigl(
              -\Delta\omega_z\,(y^2+z^2) + \Delta\omega_y\,x\,y + \Delta\omega_x\,x\,z
              \Bigr) \label{eq:dv-dr}
\end{align}

\begin{proposition}[旋转误差界]
全景图像边缘处（$|x/z|$或$|y/z|$较大），旋转标定误差被焦距放大：
\[
\|\Delta\vect{p}_i\|_{\max} \approx \|\Delta\vect{\omega}\| \cdot f_{\max} \cdot \sec^2(\text{FOV}/2)
\]
若要求拼接误差$<1$\,px，对$f=1000$\,px的相机，旋转标定精度需达到$\sim 0.06^\circ$。
\end{proposition}

\subsubsection{平移标定误差的影响}

平移误差$\Delta\vect{T}$仅影响非纯旋转情形（\cref{sec:translate-only}, \ref{sec:rt-general}）。对于平面诱导单应$\mat{H}_{21}^{\Pi}$：
\begin{equation}
\Delta\mat{H}_{21}^{\Pi} \approx \mat{K}_2\,\frac{\Delta\vect{T}_{21}\,\vect{n}\trans}{d}\,\mat{K}_1\inv
\label{eq:t-error-H}
\end{equation}

平移误差的影响与$1/d$成正比——场景越远，平移标定精度要求越低。

\begin{corollary}[标定精度需求汇总]
\centering
\small
\begin{tabular}{p{3cm}p{3.5cm}p{3.5cm}p{3cm}}
\toprule
\textbf{参数} & \textbf{影响方式} & \textbf{误差放大因子} & \textbf{典型精度需求} \\
\midrule
$f_x, f_y$ & 图像边缘放大 & $\tan(\text{FOV}/2)$ & $<0.5\%$相对误差 \\
$c_x, c_y$ & 全局平移 & $1$ & $<0.5$\,px \\
$\mat{R}$   & 图像边缘放大 & $f\cdot\sec^2(\text{FOV}/2)$ & $<0.05^\circ$ \\
$\vect{T}$  & 与$1/d$成正比 & $f/d$ & 远距离不敏感 \\
畸变系数   & 图像边缘放大 & $\propto r^3$（径向3次） & $<1\%$相对误差 \\
\bottomrule
\end{tabular}
\label{tab:calib-precision}
\end{corollary}

% --------------------------------------------------
\subsection{标定流程与LUT生成管线的集成}
\label{sec:calib-pipeline}

\begin{figure}[H]
\centering
\small
\begin{tikzpicture}[
  node distance=5mm,
  box/.style={
    draw=black!70, fill=black!3, rounded corners=3pt,
    text width=13cm, align=left, inner sep=8pt,
    font=\footnotesize
  },
  arrow/.style={->, >=stealth, thick, black!60},
  stagelabel/.style={font=\bfseries\footnotesize, inner sep=0pt, outer sep=0pt},
]

% --- 标定阶段 ---
\node[stagelabel] (label1) {标定阶段};

\node[box, below=2mm of label1] (calib) {
  \textbf{1. 单相机标定（每相机独立）} \\
  \quad 棋盘格拍摄 $\to$ 角点检测 $\to$ DLT初始 $\to$ LM优化 \\
  \quad $\to$ $\{\mat{K}_i,\ \text{distCoeffs}_i\}$ \\[3pt]
  \textbf{2. 多相机外参标定} \\
  \quad 立体标定 / 全局BA $\to$ $\{\mat{R}_{iw},\ \vect{T}_{i0}\}$ \\[3pt]
  \textbf{3. 验证：} 重投影误差 $< 0.3\ \text{px}$
};

% --- Arrow 1 ---
\draw[arrow] ([yshift=-2mm]calib.south) -- node[right, font=\footnotesize] {标定参数（一次性）} ++(0,-8mm);

% --- LUT生成阶段 ---
\node[stagelabel, below=12mm of calib.south] (label2) {LUT生成阶段（离线）};

\node[box, below=2mm of label2] (lutgen) {
  \textbf{for each} $(u_p, v_p)$ in 全景图: \\
  \quad $\vect{d}_w = \text{equirectangular}(u_p, v_p)$ \\
  \quad \textbf{for each} camera $i$: \\
  \qquad $\vect{d}_i = \mat{R}_{iw}\,\vect{d}_w$ \\
  \qquad $(x_u, y_u) = (d_{i,x}/d_{i,z},\ d_{i,y}/d_{i,z})$ \\
  \qquad $(x_d, y_d) = \mathcal{D}(x_u, y_u \mid \text{distCoeffs}_i)$ \\
  \qquad $\text{map}_x^{(i)}[v_p,u_p] = f_{x,i}\,x_d + c_{x,i}$ \\
  \qquad $\text{map}_y^{(i)}[v_p,u_p] = f_{y,i}\,y_d + c_{y,i}$ \\
  $\to$ 保存LUT文件（\texttt{float32}数组）
};

% --- Arrow 2 ---
\draw[arrow] ([yshift=-2mm]lutgen.south) -- node[right, font=\footnotesize] {预计算LUT（一次性）} ++(0,-8mm);

% --- 运行阶段 ---
\node[stagelabel, below=12mm of lutgen.south] (label3) {运行阶段（每帧）};

\node[box, below=2mm of label3] (runtime) {
  \textbf{for each} frame: \\
  \quad \textbf{for each} camera $i$: \\
  \qquad $\texttt{cv::remap}(\vect{I}_i,\ \vect{P},\ \text{map}_x^{(i)},\ \text{map}_y^{(i)},\ \texttt{INTER\_LINEAR})$ \\
  \quad $\texttt{blend}(\vect{P})$ $\to$ \textbf{输出全景图}
};

\end{tikzpicture}
\caption{标定—LUT生成—运行时拼接的完整管线}
\label{fig:pipeline}
\end{figure}

\begin{finalbox}
标定参数与LUT各章公式的对应关系：
\begin{center}
\small
\begin{tabular}{lll}
\toprule
\textbf{标定输出} & \textbf{符号} & \textbf{LUT公式引用} \\
\midrule
内参矩阵       & $\mat{K}_i$ & \cref{eq:H21}, \eqref{eq:multi-rot-main} \\
畸变系数       & $(k_1,k_2,k_3,p_1,p_2)$ & \cref{eq:one-step-main}, \eqref{eq:undistort-x} \\
相机间旋转     & $\mat{R}_{ij}$ & \cref{eq:H21} \\
相机间平移     & $\vect{T}_{ij}$ & \cref{eq:H-plane-trans}, \eqref{eq:H-plane-general} \\
世界—相机旋转 & $\mat{R}_{iw}$ & \cref{eq:multi-rot-main}, \eqref{eq:di-from-dw} \\
阵列内各相机位移 & $\vect{T}_{i0}$ & \cref{eq:Hi0-plane}, \eqref{eq:multi-depth} \\
重投影RMSE    & $\varepsilon_{\text{reproj}}$ & 误差传播：\cref{sec:calib-error-propagation} \\
\bottomrule
\end{tabular}
\end{center}
\end{finalbox}

\include{chapters/pure-rotation}
\include{chapters/pure-rotation-distortion}
\include{chapters/pure-translation}
\include{chapters/rotation-translation}
\include{chapters/lut-summary}
\include{chapters/feature-matching}

\appendix
% ============================================================
% Appendix: 补充推导、工程实现细节与符号表
% ============================================================
\section{平面诱导单应的Sherman--Morrison化简}
\label{apx:sherman-morrison}

平面诱导单应中的核心矩阵$\mat{A} = \mat{I} + \vect{t}\,\vect{n}\trans/d$（其中$\vect{t} = \vect{T}_{21}$）的逆矩阵可利用Sherman--Morrison公式化简：

\begin{equation}
\mat{A}\inv = \Bigl(\mat{I} + \frac{\vect{t}\,\vect{n}\trans}{d}\Bigr)\inv
            = \mat{I} - \frac{\vect{t}\,\vect{n}\trans}{d + \vect{n}\trans\vect{t}}
\label{eq:sherman}
\end{equation}

推广到含旋转的矩阵$\mat{B} = \mat{R} + \vect{t}\,\vect{n}\trans/d$：
\begin{equation}
\mat{B} = \mat{R}\Bigl(\mat{I} + \frac{\mat{R}\trans\vect{t}\,\vect{n}\trans}{d}\Bigr)
        \;\Longrightarrow\;
\mat{B}\inv = \Bigl(\mat{I} - \frac{\mat{R}\trans\vect{t}\,\vect{n}\trans}{d + \vect{n}\trans\mat{R}\trans\vect{t}}\Bigr)\mat{R}\trans
\label{eq:sherman-rt}
\end{equation}

这在计算后向LUT时可直接使用，无需数值求逆。

\section{LUT预计算的数值精度考虑}
\label{apx:precision}

\begin{enumerate}[leftmargin=*]
  \item \textbf{浮点精度}：LUT应使用\texttt{float32}存储。对于$4\text{K}\times 2\text{K}$的全景图，每个相机的LUT约$2 \times 4000 \times 2000 \times 4 = 64\,\text{MB}$。多相机时内存需求线性增长，可通过下采样LUT + 运行时插值来折中。
  \item \textbf{边界处理}：超出图像范围的LUT条目应标记为无效值（建议$-\text{NaN}$或$-1$），运行时通过\texttt{cv::remap}的边界模式或手动检查跳过。
  \item \textbf{畸变迭代收敛判据}：Newton迭代以$\|\mathcal{U}(x_d^{(k)}, y_d^{(k)}) - (x_u, y_u)\| < 10^{-6}$为收敛标准（对应亚像素精度）。
  \item \textbf{球面投影的极点奇异性}：在$\phi \to \pm\pi/2$（天顶/天底），$\cos\phi \to 0$使得方向向量退化为$[0, \pm 1, 0]\trans$。此时方位角$\theta$失去意义，LUT中该区域所有$\theta$值应映射到同一像素。
\end{enumerate}

\section{工程实现伪代码}
\label{apx:pseudocode}

以下给出球面全景多相机LUT生成的完整伪代码（无畸变情形）：

\begin{verbatim}
def build_lut(K_list, R_iw_list, image_sizes, pano_size):
    """
    构建多相机球面全景LUT。

    参数:
      K_list:       list of 3x3 内参矩阵 (每相机)
      R_iw_list:    list of 3x3 世界→相机旋转矩阵
      image_sizes:  list of (W, H) 每相机图像尺寸
      pano_size:    (W_p, H_p) 全景图尺寸

    返回:
      maps: list of (map_x, map_y, mask) 每相机LUT
    """
    W_p, H_p = pano_size
    f_p = W_p / (2 * pi)
    c_xp, c_yp = W_p / 2, H_p / 2

    maps = []
    for K, R_iw, (W_i, H_i) in zip(K_list, R_iw_list, image_sizes):
        map_x = np.full((H_p, W_p), np.nan, dtype=np.float32)
        map_y = np.full((H_p, W_p), np.nan, dtype=np.float32)

        for v_p in range(H_p):
            phi = (v_p - c_yp) / f_p
            cos_phi, sin_phi = cos(phi), sin(phi)

            for u_p in range(W_p):
                theta = (u_p - c_xp) / f_p

                # 球面方向向量
                d_w = np.array([
                    cos_phi * sin(theta),
                    sin_phi,
                    cos_phi * cos(theta)
                ])

                # 转换到相机坐标系
                d_i = R_iw @ d_w

                # 检查方向是否朝前 (z > 0)
                if d_i[2] <= 0:
                    continue

                # 投影到像素
                u_i = K[0,0] * d_i[0] / d_i[2] + K[0,2]
                v_i = K[1,1] * d_i[1] / d_i[2] + K[1,2]

                if 0 <= u_i < W_i and 0 <= v_i < H_i:
                    map_x[v_p, u_p] = u_i
                    map_y[v_p, u_p] = v_i

        maps.append((map_x, map_y))
    return maps
\end{verbatim}

\section{含畸变合并的LUT生成补充伪代码}
\label{apx:pseudocode-dist}

在上述伪代码的基础上，将投影步骤替换为含畸变版本：

\begin{verbatim}
# ... 在计算 d_i 之后 ...
x_u = d_i[0] / d_i[2]  # 无畸变归一化坐标
y_u = d_i[1] / d_i[2]

# 畸变施加 (迭代求解)
x_d, y_d = apply_distortion(x_u, y_u, dist_coeffs)

# 投影到畸变图像像素
u_i = K[0,0] * x_d + K[0,2]
v_i = K[1,1] * y_d + K[1,2]
# 其余相同 ...
\end{verbatim}

% --------------------------------------------------
\section*{符号表}
\label{apx:symbols}

\renewcommand{\arraystretch}{1.4}
\begin{tabularx}{\textwidth}{lX}
\toprule
\textbf{符号} & \textbf{说明} \\
\midrule
$\vect{p}_i = [u_i, v_i, 1]\trans$ & 相机$i$图像像素的齐次坐标 \\
$\mat{K}_i$ & 相机$i$内参矩阵 \\
$\vect{d}_i = \mat{K}_i\inv\vect{p}_i$ & 相机$i$归一化方向向量 \\
$\mat{R}_{ij}$ & 相机$j$坐标系到相机$i$坐标系的旋转（$\vect{X}_i = \mat{R}_{ij}\vect{X}_j$） \\
$\vect{T}_{ij}$ & 相机$j$原点在相机$i$坐标系下的平移 \\
$\mat{R}_{iw}$ & 世界坐标系到相机$i$坐标系的旋转 \\
$\vect{d}_w$ & 世界/全景坐标系下的单位方向向量（球面） \\
$(\theta, \phi)$ & 全景经纬角（方位角、仰角） \\
$(f_p, c_{x,p}, c_{y,p})$ & 全景投影参数 \\
$f_{x,i}, f_{y,i}, c_{x,i}, c_{y,i}$ & 相机$i$焦距与主点 \\
$(k_1, k_2, k_3, p_1, p_2)$ & 畸变系数（径向+切向） \\
$\mathcal{D}$ & 畸变施加映射 \\
$\mathcal{U}$ & 去畸变映射 \\
$\vect{n}$ & 平面单位法向量 \\
$d$ & 相机光心到平面的距离 \\
$\lambda_i$ & 相机$i$坐标系下空间点的深度 \\
$\mat{H}_{ij}$ & 相机$j$到相机$i$的单应矩阵（纯旋转） \\
$\mat{H}_{ij}^{\Pi}$ & 相机$j$到相机$i$的平面诱导单应矩阵 \\
$\proj^2$ & 二维射影空间 \\
$\mathbb{S}^2$ & 二维单位球面 \\
\bottomrule
\end{tabularx}


\end{document}
